\section{K-015: infinitesimal rigidity under edge doubling}

Call a support $\Sigma\subseteq\{1,\ldots,n\}^3$ \emph{free} if
no two distinct edges agree in two coordinates.  Let
$e=(i,j,k)\in\Sigma$.  We impose the following pair-deletion coverage
condition: for each mode and each coordinate in that mode, some
support edge uses that coordinate and avoids the two selected
coordinates in the other modes.  For example, for every $c$ there
is an edge $(a,b,c)\in\Sigma$ with $a\ne i$, $b\ne j$; impose the
two cyclic versions as well.

\begin{lemma}[K-015a]
Let $T=\sum_{(a,b,c)\in\Sigma}e_a\otimes e_b\otimes e_c$ have
free support, and suppose its infinitesimal product-linear
stabilizer is the two-dimensional scalar kernel.  If pair-deletion
coverage holds at $e=(i,j,k)$, replace $e$ by the four edges
\[
 (i,j,N),\quad(i,N,k),\quad(N,j,k),\quad(N,N,N),
 \qquad N=n+1.
\]
The resulting coefficient-one tensor $T'$ again has infinitesimal
stabilizer equal to the scalar kernel.
\end{lemma}

\begin{proof}
Write an infinitesimal stabilizer of $T'$ in block form
\[
 X_A=\begin{pmatrix}A&u_A\\v_A^t&\alpha\end{pmatrix},\quad
 X_B=\begin{pmatrix}B&u_B\\v_B^t&\beta\end{pmatrix},\quad
 X_C=\begin{pmatrix}C&u_C\\v_C^t&\gamma\end{pmatrix}.
\]
The coefficient of $(a,N,N)$ with $a\ne i$ in
$(X_A\otimes I\otimes I+I\otimes X_B\otimes I+
I\otimes I\otimes X_C)T'$ is $(u_A)_a$.  Thus $u_A$ is supported
at $i$; cyclically $u_B,u_C$ are supported at $j,k$.

For $a\ne i$, $b\ne j$, the coefficient at $(a,b,N)$ is
\[
                  \sum_c (v_C)_c T_{abc}.
\]
Freeness says that at most one summand is present for a fixed
$(a,b)$.  Pair-deletion coverage says that every $c$ occurs for
some such pair.  Therefore $v_C=0$.  Cyclically $v_A=v_B=0$.
Returning to $(i,N,N)$ and its cyclic counterparts now gives
$u_A=u_B=u_C=0$.  All three matrices are block diagonal.

Let $T_0=T-e_i\otimes e_j\otimes e_k$.  The component supported
on the three old spaces says that $(A,B,C)$ infinitesimally fixes
$T_0$.  The component with the third index new says
\[
 Ae_i\otimes e_j+e_i\otimes Be_j+\gamma e_i\otimes e_j=0.
\]
Projecting to either quotient by the indicated line shows
$Ae_i=ae_i$, $Be_j=be_j$, with $a+b+\gamma=0$.
The other two one-new-index components give $Ce_k=ce_k$ and
\[
 a+\beta+c=0,\qquad \alpha+b+c=0.
\]
The all-new component gives $\alpha+\beta+\gamma=0$.
Adding these equations yields $a+b+c=0$.  Hence $(A,B,C)$ also
fixes the deleted old edge, and therefore fixes $T$.  By the
hypothesis $A=aI$, $B=bI$, $C=cI$, $a+b+c=0$.  The displayed
equations then force $\alpha=a$, $\beta=b$, $\gamma=c$.
Thus the stabilizer of $T'$ is exactly the scalar kernel.
\end{proof}

\subsection*{An infinite family issuing from the negative four-dimensional component}
Use zero-based indices in the following explicit description.  Start
with the K-004 support
\[
 \Sigma_4=\{000,031,113,121,203,212,230,311,323,332\}.
\]
Its masses in this order are
\[
 (3/4,1/4,3/8,5/8,1/4,1/2,1/4,1/8,3/8,1/2).
\]
First double the edge $113$, introducing index $4$.
Thereafter double the newest diagonal edge at every step.  Write
$\Sigma_n$ for the resulting support and $T_n$ for its coefficient-one
tensor.  The $n=5$ member is type $63$ of K-010, so its Lie
stabilizer is the scalar kernel by the exact rank-$73$ certificate.

All supports remain free.  Indeed the four replacement edges agree
pairwise in at most one coordinate, and an old edge could agree
with a replacement edge in two old coordinates only if it were
the deleted edge.  Each old incidence column set remains independent:
in a relation among the enlarged columns, the three new-coordinate
rows force the coefficients of the three one-new edges to be a
common value $q$ and the all-new coefficient to be $-q$.  The old
rows then give a relation among the old columns with coefficient
$2q$ on the deleted column.  Old independence forces $q=0$ and
then all coefficients vanish.  Thus $|\Sigma_n|=3n-2$ and its
incidence rank is $3n-2$ for every $n\ge4$.

If the deleted edge has mass $r$, its replacements have masses
$r/2,r/2,r/2,1-r/2$.  This preserves every margin and strict
positivity; all masses stay below one.  In particular every support
coordinate has degree at least two.  At the newest diagonal edge
of $\Sigma_n$, pair-deletion coverage always holds for $n\ge5$.
For an old coordinate, use any surviving edge of
$\Sigma_{n-1}$ after the previously selected edge was deleted; all
coordinates remain covered because their degrees were at least two.
For the newest coordinate, use the one-new-index replacement edge,
whose other two coordinates are old.  Lemma K-015a now proves
\[
          \operatorname{Lie}\operatorname{Stab}_{GL_n^3}(T_n)
                         =\operatorname{Lie}S
                       \quad\text{for every }n\ge5.         \tag{K-015b}
\]

The positive masses give critical tensors by taking their square
roots; freeness makes all off-diagonal reduced-matrix entries zero.
Full incidence rank puts the coefficient-one tensor in their
diagonal orbit.  Thus all these tensors are polystable.
The character-lattice calculation K-006 gives diagonal stabilizer
$S\times\mu_{2^{n-1}}$, with faithful determinant character.
In particular these are actual closed tensor orbits, not just
formal supports with a large diagonal character.

For explicit spectral bookkeeping, assign the first four weights
in the three factors to be
\[
 (-1,8,-7,-4),\qquad(-1,8,-4,5),\qquad(2,-4,-1,8),
                                                               \tag{K-015c}
\]
and append the common weights
\[
              t_r=8(-2)^{r-4},\qquad 5\le r\le n.
                                                               \tag{K-015d}
\]
Take these residues modulo $M=2^{n-1}$.  Their total is
$-7+16(-2)^{n-4}\equiv-7\pmod M$, an odd number, so an odd unit
normalizes the determinant character to a chosen primitive root.
All three spectra are simple.  For $n=5$ the unique antipodal pair
in each factor consists of its old weight $8$ and the appended zero
weight.  For $n\ge6$ it consists of $t_{n-1}\equiv M/2$ and
$t_n\equiv0$.  The earlier tail weights have distinct two-adic
valuations, at least four, whereas the seed weights have valuations
at most three.  Distinct seed weights differ by less than $16$.
These observations show that there are no other coincidences or
antipodal pairs.

The fixed tensor coordinates are exactly $\Sigma_n$.  One checks
the finite seed triples directly: the nine undeleted seed edges
have sum zero, and no other seed triple does so modulo $M\ge16$.
No cross-factor seed pair has sum zero, and its only possible
nonzero sum divisible by $16$ is $16=8+8$.  With exactly one tail
coordinate, zero sum therefore occurs only for $-16+8+8=0$,
giving the three first replacement edges.  With two
tail coordinates and one seed coordinate, reduction modulo $16$
excludes zero.  For three tail coordinates, the lowest two-adic
valuation must occur at least twice.  If it occurs exactly twice,
the third weight must be the next tail weight because
$t_{r+1}=-2t_r$.  Three equal tail weights other than the newest
one have sum $3t_r\not\equiv0\pmod M$; if all three coordinates are the newest index,
their weights are all zero modulo $M$.  These give exactly the
successive triples $(r,r,r+1)$ and their three placements, and
the last diagonal edge.  There are $3n-2$ fixed coordinates.

There is also a nonzero support edge meeting exactly one of the
three paired subspaces.  At $n=5$ this follows, for example, from
the surviving edge $121$.  At $n=6$ use one of the three first
replacement edges.  For $n\ge7$ use a chain edge with two entries
at index $n-3$ and one at index $n-2$ (zero-based), the latter
belonging to the newest pair.  Consequently K-013 proves the full
stabilizer $S\times\mu_{2^{n-1}}$ for every odd $n\ge5$.
It gives a smooth quotient fixed curve of density one in each such
dimension.  The even-dimensional extension of this last step is
handled separately in K-016; the Lie statement (K-015b) itself
already holds in both parities.
